%==============================================================================
%  ملحق - تكملات: المعطيات والواجهة والأمان والاختبارات
%==============================================================================

%==============================================================================
\chapter{ملحق --- تكملات: المعطيات والواجهة والأمان والاختبارات}
\label{chap:complements}
%==============================================================================
يجمع هذا الملحق مُخرَجات تكميلية تعزّز الملف: النموذج المنطقي للمعطيات، ونماذج
الواجهة المبدئية، ومصفوفة حقوق الوصول، وسجلّ المخاطر، والامتثال لحماية المعطيات،
واستراتيجية الاختبارات.

%------------------------------------------------------------------------------
\section{النموذج المنطقي للمعطيات (MLD)}
%------------------------------------------------------------------------------
يترجم النموذج المنطقي قاموس المعطيات (الفصل 5) إلى جداول علائقية بمفاتيح أوّلية
ومفاتيح خارجية وقيود تسيير (\texttt{CHECK} على عدد العاسلات وسعة الإطارات
والإنتاجية، إضافةً إلى وحدانية زوج المقصورة/الموضع). يُميّز علاقة \emph{التموضع}
(موقع يستضيف خلايا) عن علاقة \emph{الملكية} (مزرعة تملك خلايا)، وفق القاعدة التي
تتيح لموقع استقبال خلايا من عدة مزارع.

\begin{figure}[htbp]
\centering
\includegraphics[width=0.82\linewidth]{diagramme_mld}
\caption{النموذج المنطقي العلائقي لمعطيات نظام Zümm.}
\label{fig:mld}
\end{figure}

\subsection{جدول تقابل المصطلحات}
لضمان الاتّساق بين المُخرَجات، يحمل كل كيان مهني المفهوم نفسه في ثلاثة أشكال: اسم
\textbf{قاموس المعطيات} (الفصل 5)، و\textbf{الصنف} في مخطّط UML، و\textbf{الجدول}
في النموذج المنطقي. اصطلاح كتابة الخصائص: \texttt{camelCase} في الأصناف
(\texttt{nbHausses})، \texttt{snake\_case} في الجداول (\texttt{nb\_hausses}).

\begin{xltabular}{\textwidth}{>{\bfseries}X p{4.4cm} p{4.4cm}}
\caption{تقابل الأسماء بين القاموس والأصناف والنموذج المنطقي.}
\label{tab:vocabulaire}\\
\toprule
\textbf{القاموس (مهني)} & \textbf{الصنف (UML)} & \textbf{الجدول (MLD)} \\
\midrule
\endfirsthead
\toprule
\textbf{القاموس (مهني)} & \textbf{الصنف (UML)} & \textbf{الجدول (MLD)} \\
\midrule
\endhead
مربّي النحل & \texttt{Fermier} & \texttt{FERMIER} \\
المزرعة & \texttt{Ferme} & \texttt{FERME} \\
الموقع & \texttt{Site} & \texttt{SITE} \\
الخلية & \texttt{Ruche} & \texttt{RUCHE} \\
المقصورة (جسم / عاسلة) & \texttt{Compartiment} & \texttt{COMPARTIMENT} \\
الإطار & \texttt{Cadre} & \texttt{CADRE} \\
العون & \texttt{Agent} & \texttt{AGENT} \\
البرنامج & \texttt{Planning} & \texttt{PLANNING} \\
الزيارة & \texttt{Visite} & \texttt{VISITE} \\
الصورة & (تقرير الزيارة) & \texttt{PHOTO} \\
القياس & \texttt{Mesure} & \texttt{MESURE} \\
الحصاد & \texttt{Recolte} & \texttt{RECOLTE} \\
العتبة & \texttt{Seuil} & \texttt{SEUIL} \\
\bottomrule
\end{xltabular}

%------------------------------------------------------------------------------
\section{نماذج الواجهة المبدئية (wireframes)}
%------------------------------------------------------------------------------
تُحدّد النماذج المبدئية منخفضة الدقّة التالية أرغونومية الشاشات الرئيسية قبل
التطوير. تحترم الهوية البصرية والبساطة المطلوبة.

\begin{figure}[htbp]
\centering
\includegraphics[width=0.9\linewidth]{wireframe_calendrier}
\caption{نموذج مبدئي: الرزنامة / مصفوفة الأعوان $\times$ الخلايا مع نافذة زيارة منبثقة.}
\label{fig:wf-calendrier}
\end{figure}

\begin{figure}[htbp]
\centering
\includegraphics[width=0.9\linewidth]{wireframe_tdb_production}
\caption{نموذج مبدئي: لوحة قيادة الإنتاج مع عتبة قابلة للضبط.}
\label{fig:wf-production}
\end{figure}

\begin{figure}[htbp]
\centering
\includegraphics[width=0.95\linewidth]{wireframe_rapport_visite}
\caption{نموذج مبدئي: استمارة تقرير الزيارة (مع الوضع دون اتصال والصور).}
\label{fig:wf-rapport}
\end{figure}

\begin{figure}[htbp]
\centering
\includegraphics[width=0.85\linewidth]{wireframe_ruche}
\caption{نموذج مبدئي: بطاقة خلية وتركيبها (جسم، عاسلات، إطارات).}
\label{fig:wf-ruche}
\end{figure}

%------------------------------------------------------------------------------
\section{مصفوفة حقوق الوصول (RBAC)}
%------------------------------------------------------------------------------
تحدّد المصفوفة، لكل وظيفة، الحقّ الممنوح لكل ملف. \textbf{المفتاح:} CRUD =
إنشاء / قراءة / تحديث / حذف، R = قراءة، X = تنفيذ، A = مصادقة، --- = لا وصول.

\begin{xltabular}{\textwidth}{>{\bfseries}X C{1.3cm} C{1.3cm} C{1.3cm} C{1.3cm} C{1.3cm} C{1.3cm}}
\toprule
\textbf{الوظيفة} & \textbf{المالك} & \textbf{المربّي} & \textbf{المشرف} & \textbf{المسؤول} & \textbf{المدير} & \textbf{API} \\
\midrule
\endhead
المزارع والمواقع & CRUD & R & R & R & CRUD & --- \\
الخلايا (أجسام/إطارات) & CRUD & R & R & R & CRUD & --- \\
الأعوان & R & --- & R & R & CRUD & --- \\
تخطيط زيارة & R & X & X & X & X & --- \\
المصادقة على البرنامج & --- & --- & A & --- & X & --- \\
إنجاز الزيارة/التقرير & R & X & X & --- & R & --- \\
رزنامة الأعوان $\times$ الخلايا & R & R & R & R & R & --- \\
لوحة الإنتاج & R & --- & --- & R & R & --- \\
لوحة التنبيهات / المعالجة & R & --- & --- & X & R & --- \\
الخلاصة والعائد & R & --- & --- & R & R & --- \\
ضبط العتبات/الوحدات & --- & --- & --- & R & CRUD & --- \\
التدويل & --- & --- & --- & --- & CRUD & --- \\
تصدير CSV / TXT & R & R & R & R & R & --- \\
النسخ الاحتياطي / الاستعادة & --- & --- & --- & --- & X & --- \\
واجهة كمية العسل & --- & --- & --- & --- & --- & R \\
\bottomrule
\end{xltabular}

\begin{encadre}
يجب تطبيق هذه المصفوفة من جهة الخادم (تحقّق منهجي من الحقوق في طبقة التحكم)، لا
فقط بإخفاء عناصر الواجهة.
\end{encadre}

%------------------------------------------------------------------------------
\section{سجلّ المخاطر}
%------------------------------------------------------------------------------
\begin{xltabular}{\textwidth}{>{\bfseries}p{4.6cm} C{1.8cm} C{1.6cm} X}
\toprule
\textbf{الخطر} & \textbf{الاحتمال} & \textbf{الأثر} & \textbf{إجراء التقليل} \\
\midrule
\endhead
تأخّر التطوير & متوسّط & قوي & إعطاء الأولوية لنواة CRUD، بمراحل أسبوعية (س5 إلى س12). \\
نطاق واسع جدًا & متوسّط & قوي & الاقتصار على وظائف الأولوية العالية، والباقي اختياري. \\
تعقيد التدويل العربي (RTL) & متوسّط & متوسّط & اختبار مبكر من الطرف إلى الطرف FR/EN/AR، معالجة \texttt{dir=rtl}. \\
فقدان المعطيات الميدانية & متوسّط & قوي & طابور أوامر دون اتصال + مزامنة + نسخ احتياطي. \\
ثغرات تطبيقية (حقن، XSS) & متوسّط & قوي & \texttt{PreparedStatement}، تهريب JSP، مراجعة أمنية. \\
تنبيهات كاذبة (مستشعرات) & متوسّط & متوسّط & عتبات قابلة للضبط + تحقّق بشري عبر التقرير. \\
عدم توفّر OAuth Google & ضعيف & متوسّط & تراجع إلى مصادقة محلية للعرض. \\
حجم القياسات & ضعيف & متوسّط & أرشفة وتجميع حسب الفترة (الشقّ الآلي). \\
استعمال مولّد شيفرة & ضعيف & قوي جدًا & تصميم وشيفرة يدويان، تتبّع ومراجعة الأقران. \\
\bottomrule
\end{xltabular}

%------------------------------------------------------------------------------
\section{حماية المعطيات الشخصية}
%------------------------------------------------------------------------------
يعالج النظام معطيات شخصية (أسماء وجهات اتصال مربّي النحل والأعوان) ومعطيات قد
تكون حسّاسة (\textbf{تحديد الموقع} الدقيق للمواقع). تُعتمَد المبادئ التالية،
بروح النظام العام لحماية المعطيات (RGPD):
\begin{itemize}
  \item \textbf{التقليل والغرض}: جمع المعطيات المفيدة لمتابعة تربية النحل فقط.
  \item \textbf{الأمان}: تشفير التبادلات (SSL / X.509) ووصول موثَّق، مُقرَّران سلفًا.
  \item \textbf{التحكّم في الوصول}: تطبيق صارم لمصفوفة RBAC أعلاه.
  \item \textbf{حقوق الأشخاص}: الاطّلاع والتصحيح و\textbf{المحو}، و\textbf{قابلية النقل} مضمونة بتصدير CSV / TXT.
  \item \textbf{الحفظ}: مدة احتفاظ محدَّدة للقياسات والزيارات، تُؤرشَف بعدها المعطيات أو تُخفى هويتها.
  \item \textbf{التتبّع}: تسجيل الوصول والإجراءات الحسّاسة (المصادقة على البرنامج، إغلاق موقع).
\end{itemize}

تُترجَم المبادئ أعلاه إلى الإجراءات التقنية التالية، حسب فئة المعطى. يستهدف
\textbf{التشفير عند السكون} أولًا الأعمدة الحسّاسة (جهات الاتصال، تحديد الموقع).
يفصل \textbf{إخفاء الهوية الجزئي} هوية العون عن التاريخ الزراعي عند إخفاء الهوية.

{\small
\begin{xltabular}{\textwidth}{>{\bfseries}p{3.1cm} p{3.3cm} p{2.7cm} X}
\caption{إجراءات تقنية لحماية المعطيات (بروح RGPD).}
\label{tab:rgpd}\\
\toprule
\textbf{المعطى} & \textbf{الغرض} & \textbf{الحفظ} & \textbf{الأمان / المحو} \\
\midrule
\endfirsthead
\toprule
\textbf{المعطى} & \textbf{الغرض} & \textbf{الحفظ} & \textbf{الأمان / المحو} \\
\midrule
\endhead
الهوية (الاسم، جهة الاتصال) & الحسابات والمسؤوليات & العلاقة + سنة & تشفير عند السكون، RBAC، إخفاء هوية عند الطلب \\
تحديد موقع المواقع & المتابعة، الخرائط، الطقس & حياة الموقع + سنة & تشفير عند السكون، دقّة قابلة للتخفيض، محو عند الإغلاق \\
الزيارات والتقارير & تتبّع صحي وإنتاجي & 5 سنوات & تشفير عند السكون، إخفاء هوية جزئي (فصل العون) \\
صور التفتيش & متابعة بصرية للطائفة & سنتان & تخزين مشفّر، وصول مقيَّد، حذف مع التقرير \\
قياسات المستشعرات & تحليل الأداء & خام سنة، ثم مُجمَّعة & تشفير عند السكون، تجميع / إخفاء هوية \\
سجلّات الوصول & الأمان والمساءلة & 6 إلى 12 شهرًا & وصول المدير، تطهير آلي \\
\bottomrule
\end{xltabular}}

\begin{encadre}
\textbf{إجراء المحو:} بناءً على طلب شخص معنيّ، يُنتج النظام أولًا تصديرًا (قابلية
النقل CSV / TXT)، ثم يمحو السجلّات أو يُخفي هويتها حسب الجدول أعلاه. يُحتفَظ
\textbf{بسجلّ للمعالجات} ومرجع لحماية المعطيات للمشروع.
\end{encadre}

%------------------------------------------------------------------------------
\section{استراتيجية الاختبارات}
%------------------------------------------------------------------------------
تُكمّل الاستراتيجية خطة الاختبارات الوظيفية (الفصل 10) بمستويي الوحدة والتكامل.
يجعل التصميم ضعيف الاقتران (واجهات \texttt{Repository}، خدمات مفصولة) قواعد
الأعمال قابلة للاختبار منفردةً، بحقن بدائل (mocks).

\begin{xltabular}{\textwidth}{>{\bfseries}p{2.8cm} X p{4.4cm} C{2.2cm}}
\toprule
\textbf{المستوى} & \textbf{النطاق} & \textbf{أمثلة} & \textbf{الأداة} \\
\midrule
\endhead
الوحدة & قواعد تسيير منعزلة & قيود الإطارات/العاسلات، حساب العسل (Composite)، العائد، العتبات & JUnit 5، Mockito \\
التكامل & الوصول إلى المعطيات (Repository/JDBC) & CRUD محفوظ، استعلامات الحساب & JUnit + قاعدة H2 \\
النظام & حالات وظيفية من الطرف إلى الطرف & الحالات TC01 إلى TC16 من الفصل 10 & JUnit، Selenium \\
القبول & التحقّق بالعرض & مسارات العون والمربّي والمسؤول & المناقشة \\
\bottomrule
\end{xltabular}

\begin{itemize}
  \item \textbf{التغطية المستهدَفة}: 70\,\% على الأقل من طبقة الأعمال (قواعد التسيير)، مقيسة بـ \texttt{JaCoCo}.
  \item \textbf{التتبّع}: كل متطلب مرتبط بحالة اختبار واحدة على الأقل (مصفوفة الفصل 10)، مُكمَّلة باختبارات الوحدة المقابلة.
  \item \textbf{الأتمتة}: تنفيذ الاختبارات عند كل بناء (Maven)، شرط مسبق للتسليم.
\end{itemize}

\begin{encadre}
\textbf{التآزر مع التصميم:} يتيح عكس التبعية (DIP) استبدال \texttt{RucheRepository}
في الذاكرة بتنفيذات JDBC، ما يجعل الاختبارات سريعة وحتمية ومستقلّة عن القاعدة.
\end{encadre}
